\section{Skills：技能封装}
%% ============================================================

\subsection{什么是 Skill}

当某个工作流变得可重复，就不应继续依赖长 prompt 或反复的交互。Skill 将指令（\texttt{SKILL.md} 文件）、上下文和支持逻辑封装为一个可跨 CLI、IDE 和 App 一致调用的单元。

\begin{importantbox}{Skill 设计原则}
\begin{itemize}[nosep]
    \item 每个 skill 聚焦于一个任务
    \item 从 2--3 个具体用例出发
    \item 定义清晰的输入和输出
    \item 描述（description）要说明 skill 做什么、何时使用
    \item 包含用户实际会说的触发短语
\end{itemize}
\end{importantbox}

\subsection{适用场景}

Skills 特别适合以下类型的重复性工作：

\begin{itemize}
    \item 日志分类（Log triage）
    \item Release notes 起草
    \item 按 checklist 进行 PR 审查
    \item 迁移规划（Migration planning）
    \item 遥测或事件摘要（Telemetry / incident summaries）
    \item 标准调试流程
\end{itemize}

\begin{knowledgebox}{Skill 的存储位置}
\begin{itemize}[nosep]
    \item \textbf{个人 skills}：\texttt{\$HOME/.agents/skills/}
    \item \textbf{团队共享 skills}：仓库内的 \texttt{.agents/skills/} 目录
\end{itemize}
团队共享的 skills 对新成员的 onboarding 尤其有帮助。
\end{knowledgebox}

可以使用内置的 \texttt{\$skill-creator} skill 来快速生成第一版 skill。建议在本地迭代稳定后，再打包为 plugin 进行更广泛的分享。

\begin{knowledgebox}{经验法则}
如果你发现自己反复使用同一个 prompt，或反复纠正同一个工作流，它就应该成为一个 skill。
\end{knowledgebox}

\subsection{本章小结}

Skill 是将验证过的工作流封装为可复用单元的机制。聚焦单一任务、从具体用例出发、先本地迭代后再分享。

%% ============================================================
